\subsection{Quantifiers}
\label{sec:quantifiers}
\slidesrc{FOL: 14--20}

Quantification is the central part of first-order logic. With terms and predicates alone
we still cannot say:
\begin{itemize}[nosep]
  \item all men have a mother;
  \item every natural number has a prime factor.
\end{itemize}
These sentences speak about \emph{all} objects, without naming them. We need quantifiers.

\subsubsection{The three quantifiers and their syntax}
\slidesrc{FOL: 15--18}

\begin{nbdef}{Quantifiers (slide 15)}
\begin{center}
\begin{tabular}{@{}lll@{}}
\toprule
\textbf{Symbol} & \textbf{Name} & \textbf{Read} \\
\midrule
$\forall$ & universal quantifier & ``for all \ldots'' \\
$\exists$ & existential quantifier & ``there exists \ldots'' (at least one) \\
$\exists_1$ & unique quantifier & ``there exists a unique \ldots'' (exactly one) \\
\bottomrule
\end{tabular}
\end{center}
In the Z notation of the slides a quantified formula has the form
\[ \text{quantifier}\ \ \text{signature} \bullet \text{predicate} \]
where the signature is \emph{variable : type}. Example: $\forall x : \mathit{Man} \bullet
\mathit{Mortal}(x)$.
\end{nbdef}

\begin{nbexample}{reading quantified formulas (slides 17--18)}
\small
\begin{tabularx}{\linewidth}{@{}>{\raggedright\arraybackslash}p{6.2cm} X@{}}
\toprule
\textbf{Formula} & \textbf{Meaning} \\
\midrule
$\forall x : \mathit{Man} \bullet \mathit{Mortal}(x)$ & for all $x$ of type Man, $x$ is mortal (all men are mortal) \\
$\forall x : \mathit{Man} \bullet \exists_1 y : \mathit{Woman} \bullet \mathit{MotherOf}(x, y)$ & every man has exactly one mother ($y$ is the mother of $x$) \\
$\exists m : \mathit{Monitor} \bullet \mathit{MonitorState}(m, \mathit{ready})$ & there is a monitor in the ready state \\
$\forall r : \mathit{Reactor} \bullet \exists_1 t : 100..1000 \bullet \mathit{Temp}(r) = t$ & every reactor has a temperature in the range 100 to 1000 \\
$\exists n : \N \bullet n = n * n$ & some natural number is equal to its own square (true: $0$ and $1$) \\
$\exists c : \mathit{EC} \bullet \mathit{Borders}(c, \mathit{Albania})$ & some EC country borders Albania (Greece) \\
$\forall m, n : \mathit{Person} \bullet \neg \mathit{Superior}(m, n)$ & no person is superior to another \\
$\forall m : \mathit{Person} \bullet \neg \exists n : \mathit{Person} \bullet \mathit{Superior}(m, n)$ & ditto (same meaning) \\
\bottomrule
\end{tabularx}
\end{nbexample}

\paragraph{Without types.} Most AI texts (and Prolog) use untyped FOL: there is one big
domain of objects, and types become ordinary predicates. The typed formulas translate
with a fixed pattern that is worth memorising.

\begin{nbfact}{Typed to untyped quantifiers}
\notinslides
\begin{align*}
\forall x : T \bullet P(x) \quad &\text{becomes}\quad \forall x\,\big(T(x) \limp P(x)\big) \\
\exists x : T \bullet P(x) \quad &\text{becomes}\quad \exists x\,\big(T(x) \land P(x)\big)
\end{align*}
\textbf{$\forall$ goes with $\limp$, $\exists$ goes with $\land$.}
\end{nbfact}

\begin{nbwarning}[The two classic translation mistakes]
\begin{itemize}[nosep]
  \item ``All men are mortal'' is $\forall x\,(\mathit{Man}(x) \limp \mathit{Mortal}(x))$,
        \textbf{not} $\forall x\,(\mathit{Man}(x) \land \mathit{Mortal}(x))$: the latter
        says that \emph{everything} (every chair, every number) is a man and is mortal.
  \item ``Some monitor is ready'' is $\exists x\,(\mathit{Monitor}(x) \land
        \mathit{Ready}(x))$, \textbf{not} $\exists x\,(\mathit{Monitor}(x) \limp
        \mathit{Ready}(x))$: the latter is true as soon as there exists any object that is
        not a monitor (false antecedent!), even if no monitor is ready.
\end{itemize}
\end{nbwarning}

\begin{nbwarning}[The order of quantifiers matters]
\notinslides
$\forall x\,\exists y\,\mathit{MotherOf}(x, y)$: everybody has a mother (possibly a
different one for each $x$): true. \\
$\exists y\,\forall x\,\mathit{MotherOf}(x, y)$: there is one $y$ who is the mother of
everybody: false. Swapping $\forall$ and $\exists$ changes the meaning; swapping two
$\forall$ (or two $\exists$) does not.
\end{nbwarning}

\subsubsection{Free and bound variables}
\notinslides

A quantifier $\forall x$ or $\exists x$ \textbf{binds} the variable $x$ inside its
\textbf{scope} (the formula it is applied to). An occurrence of a variable that is not
bound by any quantifier is \textbf{free}. In
$\forall x\,(\mathit{gt}(x, y) \limp \mathit{gt}(x, 0))$, $x$ is bound and $y$ is free.
A formula with no free variables is a \textbf{sentence} (or \emph{closed} formula): only
sentences are simply true or false in an interpretation; a formula with free variables,
like $\mathit{gt}(x, 3)$, needs a value for its free variables (slide 13). This
distinction comes back in the proof rules of Section~\ref{sec:proof-theory-fol}.

\subsubsection{Quantifiers as big conjunctions and disjunctions}
\slidesrc{FOL: 19--20}

\begin{nbfact}{Quantifiers over a finite set (slide 19)}
Universal quantification is like conjunction, existential quantification like
disjunction:
\begin{align*}
\forall n : \{2, 4, 6\} \bullet \mathit{Even}(n) \quad &\text{is the same as}\quad
  \mathit{Even}(2) \land \mathit{Even}(4) \land \mathit{Even}(6) \\
\exists n : \{7, 8, 9\} \bullet \mathit{Prime}(n) \quad &\text{is the same as}\quad
  \mathit{Prime}(7) \lor \mathit{Prime}(8) \lor \mathit{Prime}(9)
\end{align*}
\end{nbfact}

\begin{nbfact}{Duality of $\forall$ and $\exists$ (slide 20)}
\begin{align*}
\forall x : T \bullet P(x) \quad &\liff\quad \neg\,\exists x : T \bullet \neg P(x) \\
\exists x : T \bullet P(x) \quad &\liff\quad \neg\,\forall x : T \bullet \neg P(x)
\end{align*}
Saying that everything has some property is the same as saying that there is nothing that
does not have it. Saying that something has the property is the same as saying that it is
not the case that everything lacks it.
\end{nbfact}

\begin{nbintuition}[Duality is De Morgan in disguise]
Over the finite set $\{a, b, c\}$:
\begin{align*}
\neg\,\forall x\, P(x)
  &= \neg\big(P(a) \land P(b) \land P(c)\big) && \reason{$\forall$ = big $\land$} \\
  &\equiv \neg P(a) \lor \neg P(b) \lor \neg P(c) && \reason{De Morgan} \\
  &= \exists x\, \neg P(x) && \reason{$\exists$ = big $\lor$}
\end{align*}
Negating both sides gives the first duality law. Practical reading: to refute ``all
monitors are ready'' you need one monitor that is \emph{not} ready (a counterexample).
\end{nbintuition}

\begin{nbfact}{The unique quantifier, defined with the other two}
\notinslides
$\exists_1 x\, P(x) \;\equiv\; \exists x\,\big(P(x) \land \forall y\,(P(y) \limp y = x)\big)$:
there is an $x$ with the property, and any $y$ with the property is that same $x$.
\end{nbfact}

\begin{nbexercise}{Q.1}{Translate into FOL (added)}
Use the predicates $\mathit{Student}(x)$, $\mathit{Course}(x)$, $\mathit{Attends}(x, y)$,
$\mathit{Passed}(x, y)$.
\begin{enumerate}[nosep]
  \item Every student attends some course.
  \item Some course is attended by every student.
  \item No student passed a course without attending it.
  \item Not every student passed SDAI.
\end{enumerate}
\end{nbexercise}
\begin{nbsolution}
\begin{enumerate}[nosep]
  \item $\forall x\,\big(\mathit{Student}(x) \limp \exists y\,(\mathit{Course}(y) \land
        \mathit{Attends}(x, y))\big)$
  \item $\exists y\,\big(\mathit{Course}(y) \land \forall x\,(\mathit{Student}(x) \limp
        \mathit{Attends}(x, y))\big)$ \ (stronger than 1: the order of quantifiers changed)
  \item $\neg \exists x\, \exists y\, \big(\mathit{Student}(x) \land \mathit{Passed}(x, y)
        \land \neg \mathit{Attends}(x, y)\big)$, equivalently
        $\forall x\, \forall y\, \big(\mathit{Student}(x) \land \mathit{Passed}(x, y)
        \limp \mathit{Attends}(x, y)\big)$ (by duality and De Morgan).
  \item $\neg \forall x\,(\mathit{Student}(x) \limp \mathit{Passed}(x, \mathit{sdai}))$,
        equivalently $\exists x\,(\mathit{Student}(x) \land \neg \mathit{Passed}(x,
        \mathit{sdai}))$: some student did not pass SDAI.
\end{enumerate}
\textbf{Lesson}: always check the pattern $\forall \ldots \limp$, $\exists \ldots \land$.
\end{nbsolution}
